% ----------------------------------------------------------
% Abstract for talk at 26th British Combinatorial Conference
% ----------------------------------------------------------

% ----------------------------------------------------------
% PLEASE DO NOT MODIFY THIS HEADER SECTION

\documentclass[12pt,a4paper]{article}
\newcommand{\TalkTimeLocation}{Xxxday 00:00 TC30x}

\pagestyle{empty}

\frenchspacing \setlength{\parskip}{9pt plus 3pt minus 1pt}
\setlength{\parindent}{0pt}

\setlength{\textwidth}{6.27in}
\setlength{\textheight}{9.69in}
\setlength{\topmargin}{0pt}
\setlength{\headsep}{0pt}
\setlength{\headheight}{0pt}
\setlength{\oddsidemargin}{0pt}
\setlength{\evensidemargin}{0pt}

\usepackage{amsmath}
\usepackage{amssymb}
\usepackage{amsthm}
\usepackage{tikz}

\newtheorem{theorem}{Theorem}
\newtheorem{corollary}[theorem]{Corollary}
\newtheorem{lemma}[theorem]{Lemma}
\newtheorem{proposition}[theorem]{Proposition}
\newtheorem{conjecture}[theorem]{Conjecture}
\newtheorem{question}[theorem]{Question}

\newcommand{\ifEmpty}[3]{\setbox0=\hbox{#1\unskip}\ifdim\wd0=0pt{#2}\else{#3}\fi}
\newcommand{\nothingIfEmpty}[2]{\ifEmpty{#1}{}{#2}}
\newcommand{\showIfNotEmpty}[1]{\nothingIfEmpty{#1}{#1}}
\newcommand{\showTalkJointWork}[1]{\nothingIfEmpty{#1}{(This talk is based on joint work with #1.)}}
\newcommand{\warn}[1]{\color{red}\large\textbf{\textsf{#1}}}
\renewcommand{\refname}{}
\newenvironment{talkBibliog}{\vspace*{-24pt}\begin{thebibliography}{9}\footnotesize}{\end{thebibliography}}

% PLEASE DO NOT MODIFY ABOVE HERE

% ----------------------------------------------------------
% PLEASE SPECIFY THE TALK INFORMATION HERE

\newcommand{\TalkTitle}{Spectral bounds for graph parameters -- Part II} %TITLE of the talk}
\newcommand{\TalkPresenter}{Pawel Wocjan} %NAME of the speaker}
\newcommand{\TalkEmail}{wocjan@cs.ucf.edu} %my.email@uni.ac.uk}
\newcommand{\TalkAffiliation}{University of Central Florida} %University of the speaker}
\newcommand{\TalkJointWork}{Clive Elphick} %A. Coauthor}     % leave blank if not joint work
\newcommand{\TalkMSCCodes}{05C15, 05C50} %MSC codes here}                   % REQUIRED - see www.ams.org/msc/msc.html

% ----------------------------------------------------------
% PERSONAL MACROS MAY BE PLACED HERE

% \newcommand{\iso}{\cong}

\newcommand{\C}{\mathbb{C}}

% ----------------------------------------------------------
% PLEASE DO NOT MODIFY THIS SECTION

\begin{document}
\hfill \textsf{\TalkTimeLocation}
\begin{center}
 {\Large\textsc{\TalkTitle}}

 \textbf{\TalkPresenter}

 \showIfNotEmpty{\texttt{\TalkEmail}}

 \showIfNotEmpty{\TalkAffiliation}

 \showTalkJointWork{\TalkJointWork}

 {\small{MSC2000: \ifEmpty{\TalkMSCCodes}{\warn{An MSC code is required!}}{\TalkMSCCodes}}}

 \bigskip\bigskip
\end{center}

% ----------------------------------------------------------
% PLEASE INSERT YOUR ABSTRACT HERE

This talk is the second of a pair of consecutive talks by Clive Elphick and Pawel Wocjan on their joint work on spectral bounds for graph parameters, that lie between the clique number $\omega(G)$, and the chromatic number, $\chi(G)$. The second talk will provide more technical details. 

Let $[c]=\{0,\ldots,c-1\}$ for some positive integer $c$ and let $I_d$ and $0_d$ denote the identity and zero matrices acting on $\C^d$. 
A quantum $c$-coloring of an undirected graph $G=(V,E)$ without loops and without multiple edges is a collection of orthogonal projectors 
\[
\{ P_{v,k} : v \in V, k \in [c] \}
\]
acting on $\C^d$ for some $d>0$ such that
\begin{itemize}
\item for all vertices $v\in V$
\[
\sum_{k\in [c]} P_{v,k} = I_d \quad \quad \mbox{completeness}
\]
\item for all edges $vw\in E$ and for all $k\in [c]$
\[
P_{v,k} P_{w,k} = 0_d \quad \quad \mbox{orthogonality}
\]
\end{itemize}
The quantum chromatic number, $\chi_q$, is the smallest $c$ for which the graph $G$ admits a quantum $c$-coloring. The classical chromatic number is a special case when $d=1$.

We will establish that the existence of a quantum $c$-coloring implies the existence of a unitary matrix $U$ acting on $\C^n \otimes \C^d$ such that
\[
\sum_{\ell=1}^{c-1} U^\ell (A\otimes I_d) (U^*)^\ell = -A \otimes I_d,
\]
where $A$ is the adjacency matrix of $G$ and $n$ its number of vertices.  

We will show that applying the majorization result for maximal eigenvalues of Hermitian matrices $X$ and $Y$, 
$\mu_{\max}(X) + \mu_{\max}(Y) \ge \mu_{\max}(X+Y)$, to the above matrix equality yields the Hoffman bound for the quantum chromatic number: $\chi_q(G) \ge 1 + \mu_{\max(A)}/|\mu_{\min}(A)|$.

Similarly, we will show that applying the rank inequality for positive semidefinite matrices $X$ and $Y$ with $X-Y$ positive semidefinite, $\mathrm{rank}(X)\ge\mathrm{rank}(Y)$, to the above matrix inequality yields the inertial lower bound on the quantum chromatic number: $\chi_q(G)\ge 1 + \max\{n^+(A)/n^-(A), n^-(A) / n^+(A)\}$, where $n^\pm(A)$ denote the number of positive and negative eigenvalues of $A$.

%Please do the following:
%\begin{enumerate}
%  \item Save this file as \texttt{\emph{YourFirstName}-\emph{YourLastName}.tex}, using your name.
%  \item Add the talk information as required. You must specify at least one MSC code; see \texttt{www.ams.org/msc/msc.html}.
%  \item Insert your abstract. Graphics may be included using TikZ. Theorems, etc., may be included using the appropriate environments. You may also include references. The complete abstract must fit on a single page. Do not change the font size or the margins.
%  \item If you haven't already done so, register for the conference by visiting BCC's webpage.
%  \item Email both the .tex file and the PDF to \texttt{bcc2019@contacts.bham.ac.uk} by 31st~May 2019.
%\end{enumerate}
%Thanks!





% ----------------------------------------------------------
% PLEASE ADD ANY REFERENCES HERE




% ----------------------------------------------------------
\end{document}
% ----------------------------------------------------------
